\section{Dustcasting}
\label{sec:dustcasting}

Once the party reaches their dustcasting target, read the following to the player doing the dustcasting:

\quoteBox{
	You place your palm with the \lootDustcaster{} to the cold stone. Suddenly you have a feeling that everything else is gone - there is only you and the stone and the fabrial whispering eagerly: "Make... it... gone...!"
}

A character using \lootDustcaster{} must do a \skillCheck{10}{Division} or \skillCheck{20}{\skillDiscipline{}} three times. After the 3rd test, the dustcasting is complete and their objective and \lootDustcaster{} itself are annihilated. However, based on the number of failures, additional effects take place:

\scriptedOption{1+ Failures}{
	The PC looses the tips of their fingers on the dustcasting hand. They take 1d6 Spirit damage.
}

\scriptedOption{2+ Failures}{
	Whole dustcasting palm is gone. The PC takes additional 1d8 Spirit damage.
}

\scriptedOption{3 Failures}{
	Entire arm of the PC is turned to dust. Its owner takes additional 1d10 Spirit damage.
}

\subsubsection{Dropping to 0 Health}

If the dustcasting failures reduce a character to 0 Health, that character's entire body turns to dust and the PC is gone. That doesn't stop the dustcasting.

\subsubsection{Unable to Concentrate}

If the dustcasting character takes damage (eg after guard's attack, see \fullref{sec:dustcastingBattle}), their next test for the dustcasting will gain disadvantage.

\reasonBox{
	Please note, "Dustcasting" is not an official term from the \bookSA{} lore. I based it on "soulcasting" term and modified it to indicate \npcDustmother{}, the \npcDustmotherTitle{}'s influence. I feel it fits, but it's not official.
}

\corruptedBox{
	Using \lootDustcaster{}, an artifact of \npcDustmother{}, the \npcDustmotherTitle{}, can expedite the progression of an Enlightened Dustbringer \seeCorruptedRules{}.
}



\subsection{Guards! Guards!}
\label{sec:dustcastingBattle}

If the party failed to sneak to their target (see \fullref{sec:dustcastingSneaking}), they will be attacked by the guards. The PC performing the dustcasting cannot participate in the battle in any way and they must be defended by other characters.

During the fight, each dustcasting test roll takes a turn to complete. For three rolls, the fight will last for three turns. Once the dustcasting is complete, the party can safely flee the battle.

\scriptedOption{Singer Guards}{
	If the dustcasting takes place in the Singer Camp, the characters are attacked by \enemySingerDireform{} and four \enemySingerNimbleform{}s. At beginning of each round, additional \enemySingerWarform{} joins the fight.
}

\scriptedOption{Wall Guard}{
	If the dustcasting takes place at the Kholinar walls, the characters will be attacked by six \enemyWallGuard{}s. At beginning of each round, additional two \enemyWallGuard{}s will join the fight.
}


\subsection{Aftermath}

Once the dustcasting is completed (even if a character performing it dies in a process), read the following:

\quoteBox{
	In an instant the great stony surface disappears in a puff of smoke. Huge gray cloud suddenly erupts many feet in all directions. You are sure it must be visible from far, far away. \\ \\
	But there is no time to ponder - you must flee! Now!
}

The dustcasting is complete, but every guard is now aware of the PCs' presence. They must run. Proceed to \fullref{sec:dustcastingFlee}.